\section{Appendix B: Complete-record block experiments}

The lower-bound analysis uses pairs of priors over fixed schedules whose induced experiments retain the complete observed record. The construction separates labeled sources from fixed groups of recipients, allowing the response likelihood to account jointly for retained arrows, treatment assignments, and repeated recipient outcomes. Its auxiliary parameters are the block count \leanref{sym:B}{\ensuremath{B}} and the source count \leanref{sym:m}{\ensuremath{m=Bd}}. The following definition fixes the recipient groups and the population capacity required by the construction.

\begin{definitionv}{synth_4}[Fixed recipient blocks]
For integers \(n\ge4\) and \(B,d\ge1\) satisfying \(2Bd\le n\), set \(m=Bd\) and \(V=\{1,\ldots,n\}\). The source units are \(\{1,\ldots,m\}\). Define the fixed recipient blocks by
\[
C_\ell=\{m+(\ell-1)d+1,\ldots,m+\ell d\},
\qquad \ell\in\{1,\ldots,B\}.
\]
Each \(C_\ell\) contains \(d\) units, and these pairwise disjoint blocks partition \(\{m+1,\ldots,2m\}\). The padding units are \(\{i\in V:i>2m\}\). Thus the sources, recipient blocks, and padding units partition \(V\). The recipient blocks remain fixed as the ordered source partition \(S=(S_1,\ldots,S_B)\) varies; \(S_\ell\) is the source group assigned to \(C_\ell\) in the block-prior construction.
\end{definitionv}

The recipient block \leanref{sym:Cblocks}{\ensuremath{C_\ell}}, indexed by \leanref{sym:ell}{\ensuremath{\ell}}, remains fixed while the ordered source partition \leanref{sym:S}{\ensuremath{S}} varies. This distinction makes the repeated-response groups observable from recipient labels while leaving source membership partially revealed by the audit. We next specify the partition domain and the baseline density used to randomize fixed schedules.



The baseline vector \leanref{sym:U}{\ensuremath{U}} has density determined by \leanref{sym:f}{\ensuremath{f}}, with real density argument \leanref{sym:w}{\ensuremath{w}}, as specified in \cref{obj:def:block-prior}. The sign \leanref{sym:sigma}{\ensuremath{\sigma}} selects the direction of the spillover effect, and the amplitude \leanref{sym:h}{\ensuremath{h}} controls its magnitude. Each baseline is shared by a recipient block and is drawn before the experimental assignment. Thus the baseline randomization defines a prior over schedules, with each realized schedule held fixed during its experiment.

To describe the observable information entering the likelihood, consider a retained graph compatible with the block layout. A retained outgoing arrow identifies its source's recipient block, and the assignment vector records the treatment of every labeled source. The following definition collects the resulting capacities, sign sums, and count-constrained likelihood notation.



The remaining capacities \leanref{sym:urow}{\ensuremath{u_\ell}} and revealed sign sums \leanref{sym:A}{\ensuremath{A_\ell}} in \cref{obj:lem:block-law} are functions of the labeled retained graph and assignments. The same record determines the total remaining capacity \leanref{sym:M}{\ensuremath{M}} and the global hidden treated-source count \leanref{sym:K}{\ensuremath{K}}. The block-response vector \leanref{sym:y}{\ensuremath{y}} records one outcome per fixed recipient group, preserving the repeated observations through the known groups. In the likelihood expression \leanref{sym:density}{\ensuremath{\lambda_{\sigma,h}(y\mid H,Z)}} defined in \cref{obj:lem:block-law}, the indeterminate \leanref{sym:x}{\ensuremath{x}} tracks treated counts, while \leanref{sym:k}{\ensuremath{k}} indexes a block's possible hidden treated count.

These definitions organize the information used by the complete-record experiment. The prior now specifies how the source partition and baselines jointly generate its fixed graph and coefficients.

\begin{definitionv}{def:block-prior}[Block schedule prior]
The block prior \(\Pi_{\sigma,h}\), indexed by the sign \(\sigma\in\{-1,1\}\) and amplitude \(h\in\mathbb R\), is the distribution over fixed schedules \(\theta=(G,a,t,b)\) induced by the following construction. With \(B\) blocks and \(m=Bd\) sources, draw an ordered partition \(S=(S_1,\ldots,S_B)\) uniformly over partitions of the sources satisfying \(|S_\ell|=d\). Independently of \(S\), draw independent block baselines \(U_1,\ldots,U_B\) with density
\[
f(w)=
\begin{cases}
4\cos^2(2\pi w),& |w|\le 1/4,\\
0,& |w|>1/4.
\end{cases}
\]
Using the recipient blocks \(C_\ell\) of the block construction, set
\[
G=\bigcup_{\ell=1}^{B}\{(j,i):j\in S_\ell,\ i\in C_\ell\},
\qquad
a_i=\sum_{\ell=1}^{B}\mathbf1\{i\in C_\ell\}
       \left(U_\ell-\frac{\sigma h}{2}\right),
\qquad
t_i=0,
\qquad
b_{ij}=\frac{\sigma h}{d}.
\]
Spillover coefficients are read on graph arrows. Source and padding rows have zero outcomes. The baselines are drawn before treatment randomization and are fixed components of each schedule.

The corresponding complete-record mixture law \(P_{\sigma,h}\) is obtained by drawing the partition and baselines, fixing the resulting schedule, and then independently drawing treatment assignments and audit marks from the canonical product design with retention parameter \(q\). For every measurable set \(A\) of complete observed records \(O=(H,Z,Y(Z))\),
\[
P_{\sigma,h}(A)
=
\int
\Pr\!\left\{(H,Z,Y(Z))\in A\mid\theta\right\}
\,\Pi_{\sigma,h}(d\theta),
\]
where the conditional probability uses that assignment-audit design.
\end{definitionv}

The schedule prior \leanref{sym:Pi}{\ensuremath{\Pi_{\sigma,h}}} and complete-record mixture law \leanref{sym:Psign}{\ensuremath{P_{\sigma,h}}} in \cref{obj:def:block-prior} distinguish the two sources of randomization. Partition and baseline draws select a schedule; treatment and audit draws then produce an observation conditional on that schedule. Within a recipient block, the baseline centering gives the realized response
\[
Y_i(Z)
=
U_\ell+\frac{\sigma h}{2d}\sum_{j\in S_\ell}(2Z_j-1),
\qquad i\in C_\ell.
\]
Consequently, the same block response appears at each of its \(d\) recipients. The testing pair uses the following explicit amplitude to relate this construction to the estimation scale.

\begin{definitionv}{def:testing-handle}[Testing amplitude $h_\circ(B,d,q)$]
The testing amplitude is
\[
h_\circ(B,d,q)=\frac{1}{100\pi}
\begin{cases}
\min\left\{1,\dfrac{d}{\sqrt{mp}}\right\},&p>0,\\
1,&p=0,
\end{cases}
\]
where the source count \(m\) and association-revelation probability \(p\) are
\[
m=Bd,
\qquad
p=1-(1-q)^d.
\]
The parameters satisfy
\begin{itemize}
\item \textbf{(Block count.)} \(B\) is an integer with \(B\ge1\).
\item \textbf{(Degree.)} \(d\) is an integer with \(d\ge1\).
\item \textbf{(Population capacity.)} \(2m\le n\).
\end{itemize}
The testing handle consists of the opposite-sign block priors and their complete original-record mixture laws:
\[
\bigl(\Pi_{+1,h_\circ(B,d,q)},\Pi_{-1,h_\circ(B,d,q)}\bigr),
\qquad
\bigl(P_{+1,h_\circ(B,d,q)},P_{-1,h_\circ(B,d,q)}\bigr).
\]
The record includes detailed retained arrows and source labels in the retained graph \(H\), every coordinate of the assignment vector \(Z\), all noiseless realized outcomes \(Y(Z)\), repeated-outcome recipient grouping, remaining block capacities \(u_\ell\), and the observed global hidden treated-source count \(K\). Its conditional response density \(\lambda_{\sigma,h}\) uses coefficient extraction constrained by \(K\), including the convention for zero total remaining capacity \(M=0\).

Support, causal target separation, the uniform nonlocal testing bound at \(h_\circ\), and the all-degree Walsh contraction are proved properties of this construction. Together they yield the risk scale \(F\), its equivalent positive-retention elbow expression, the clipped-or-zero rule \(T^\star_{n,d,q}\), and the two consistency criteria
\[
\frac{d_n^2}{n}\longrightarrow0,
\qquad
\frac{d_n}{nq_n}\longrightarrow0,
\]
for admissible population sequences, with the second ratio assigned \(+\infty\) at zero retention.
\end{definitionv}

The amplitude \(h_\circ(B,d,q)\) in \cref{obj:def:testing-handle} lies within the support range \(0\le h\le1/4\). At positive association-revelation probability it combines a bounded amplitude with the scale \(d/\sqrt{mp}\); at zero revelation it uses the specified constant amplitude. The support and exact likelihood are supplied by the next result, while the contraction and testing bounds appear in \cref{obj:lem:hidden-allocation-contraction,obj:thm:block-testing-scale,obj:thm:nonlocal-testing-answer}.

\begin{lemmav}{lem:block-law}[Block mixture and reconstruction]
Let \(n,B,d\) be integers, let \(\sigma\in\{-1,1\}\) be the prior sign, let \(h\) be its amplitude, and let \(q\) be the edge-retention probability. Suppose:
\begin{itemize}
\item \textbf{(Population and blocks.)} \(n\ge4\), \(B\ge1\), \(d\ge1\), and \(2Bd\le n\).
\item \textbf{(Parameter ranges.)} \(0\le h\le1/4\) and \(0\le q\le1\).
\item \textbf{(Experimental design.)} The joint law of the assignment vector \(Z\) and audit-mark array \(W\) satisfies \cref{obj:ass:assignment,obj:ass:audit,obj:ass:design-independence}.
\end{itemize}

Set \(m=Bd\), the number of sources. Partition the \(m\) labeled sources uniformly into ordered blocks \(S_1,\ldots,S_B\) of size \(d\), independently of block baselines \(U_1,\ldots,U_B\), which are independent with density
\[
f(u)=
\begin{cases}
4\cos^2(2\pi u),& |u|\le1/4,\\
0,& |u|>1/4.
\end{cases}
\]
For the fixed recipient blocks \(C_1,\ldots,C_B\), each of size \(d\), include every arrow from \(S_\ell\) to \(C_\ell\). Set
\[
a_i=U_\ell-\frac{\sigma h}{2}\quad(i\in C_\ell),
\qquad t_i=0,
\qquad b_{ij}=\frac{\sigma h}{d}
\]
on these arrows, and give source and padding units zero outcomes. Let \(\Pi_{\sigma,h}\) be the resulting schedule law and \(P_{\sigma,h}\) its complete-record mixture law. Then, for \(\Pi_{\sigma,h}\)-almost every schedule \(\theta\),
\[
\theta\in\mathcal M_{n,d},
\qquad
\tau(\theta)=\frac{\sigma mh}{n},
\]
where \(\mathcal M_{n,d}\) is the schedule class in \cref{obj:def:model} and \(\tau(\theta)\) is the population all-treated-versus-all-control contrast.

Under the actual retained-graph marginal, the indicators that each labeled source has at least one retained outgoing arrow are independent Bernoulli variables with common probability
\[
p=1-(1-q)^d.
\]

Let \(H\) be the detailed labeled retained graph. For each block, define its remaining source capacity and revealed source-sign sum by
\[
u_\ell=d-\bigl|\{j\in S_\ell:j\text{ has a retained outgoing arrow}\}\bigr|,
\qquad
A_\ell=\sum_{\substack{j\in S_\ell\\j\text{ has a retained outgoing arrow}}}(2Z_j-1).
\]
Define the total remaining capacity and hidden treated-source count by
\[
M=\sum_{\ell=1}^B u_\ell,
\qquad
K=\sum_{\substack{j\text{ a source}\\j\text{ has no retained outgoing arrow}}}Z_j.
\]
For almost every \((H,Z)\) under its actual joint marginal, the conditional density of the distinct block-response vector \(y=(y_1,\ldots,y_B)\) is
\[
\lambda_{\sigma,h}(y\mid H,Z)
=
\frac{
[x^K]\displaystyle\prod_{\ell=1}^B
\left\{
\sum_{k=0}^{u_\ell}
\binom{u_\ell}{k}
f\!\left(y_\ell-\frac{\sigma h(A_\ell+2k-u_\ell)}{2d}\right)x^k
\right\}
}{
\binom{M}{K}
},
\]
where \([x^K]\) extracts the coefficient of \(x^K\). The law \(P_{\sigma,h}\) equals the law obtained by drawing \((H,Z)\) from that actual joint marginal, drawing \(y\) from this conditional density, copying \(y_\ell\) to every recipient in \(C_\ell\), and assigning zero outcomes to source and padding units. Its \((H,Z)\) marginal equals the actual joint marginal, including detailed retained-edge subsets, full source labels, and every treatment coordinate. For almost every \((H,Z)\) under this marginal,
\[
M=0\ \Longrightarrow\ K=0\ \text{ and }\ \binom{M}{K}=1.
\]
\end{lemmav}

\Cref{obj:lem:block-law} establishes support in the schedule class and a constant causal target under each prior. The construction's nonzero degrees are \(d\), and each recipient's absolute coefficient mass is bounded by
\[
\frac14+\frac h2+h\le\frac58
\qquad(0\le h\le1/4).
\]
Each of the \(m\) recipients has all-treated-versus-all-control contrast \(\sigma h\), giving the target \(\sigma mh/n\). The opposite-sign priors therefore have target separation \(2mh/n\), including at the amplitude specified in \cref{obj:def:testing-handle}. These are the support and separation properties used by the risk reduction in \cref{obj:lem:full-record-two-prior-risk}.

The conditional likelihood expresses how the observed assignments constrain the remaining source allocation. For a block receiving \(k\) hidden treated sources, the hidden sign sum is \(2k-u_\ell\); together with the revealed sum \(A_\ell\), it determines the translation of the baseline density. The coefficient extraction admits precisely those blockwise counts whose sum equals the observed global count \(K\). Its binomial normalization accounts for allocating that fixed treated count across the remaining capacities. At zero remaining capacity, the likelihood uses the stated \(K=0\) convention.

Finally, the reconstruction clause in \cref{obj:lem:block-law} preserves the complete observation channel. Detailed audit subsets and retained-edge labels enter through the actual graph marginal, every treatment coordinate remains in the joint marginal, and each generated block response is copied to its full recipient group. Although the conditional response formula summarizes the graph through observable capacities and sign sums, the reconstructed law retains the detailed graph itself. The resulting testing experiment therefore accommodates estimation procedures that use retained labels and repeated noiseless responses to infer source associations.
